\documentclass[letterpaper]{article}
\usepackage{aaai}
\usepackage{times}
\usepackage{helvet}
\usepackage{courier}
\usepackage{booktabs}
\usepackage{microtype}
\usepackage{amsmath}
\usepackage{amssymb}
\usepackage{graphicx}
\usepackage{float}
\frenchspacing
\setlength{\pdfpagewidth}{8.5in}
\setlength{\pdfpageheight}{11in}
\pdfinfo{
/Title (Supplementary Material: bioMoR)
/Author (Anonymous Submission)
}
\setcounter{secnumdepth}{1}

\title{Supplementary Material\\
bioMoR: Biology-Guided Mixture-of-Recursions for Efficient and Interpretable Genomic Learning}
\author{Anonymous Submission}

\begin{document}
\maketitle
\appendix

\noindent This supplement collects the additional tables and derivations referenced in the
main paper. Unless stated otherwise, the auxiliary empirical tables here use a $3$-seed
held-out protocol; the $5$-fold cross-validation results reported in the main paper are the
authoritative numbers, and we include these analyses for completeness and transparency.

% Training-dynamics section removed; the Baron loss/F1 figure now lives in the main paper.

\section{Model-Size Scaling}
\label{app:scaling}
\begin{table}[t]
\centering
\IfFileExists{cv5_scaling_table.tex}{\resizebox{\columnwidth}{!}{%
\begin{tabular}{c ccc ccc}
\toprule
 & \multicolumn{3}{c}{Single-cell} & \multicolumn{3}{c}{Multi-omics (P-NET)} \\
\cmidrule(lr){2-4}\cmidrule(lr){5-7}
$d_{\text{model}}$ & Vanilla & MoR-gen. & bioMoR & Vanilla & MoR-gen. & bioMoR \\
\midrule
96 & 63.6{\scriptsize$\pm$5.9} & 63.4{\scriptsize$\pm$4.1} & 74.8{\scriptsize$\pm$2.0} & 52.1{\scriptsize$\pm$5.1} & 55.6{\scriptsize$\pm$9.0} & 56.2{\scriptsize$\pm$4.3} \\
136 & 65.1{\scriptsize$\pm$5.0} & 62.4{\scriptsize$\pm$5.6} & 75.8{\scriptsize$\pm$3.8} & 51.2{\scriptsize$\pm$7.7} & 66.0{\scriptsize$\pm$9.9} & 53.5{\scriptsize$\pm$5.6} \\
192 & 61.8{\scriptsize$\pm$5.1} & 62.8{\scriptsize$\pm$4.3} & 76.2{\scriptsize$\pm$3.3} & 42.5{\scriptsize$\pm$6.1} & \emph{\scriptsize run\ldots} & 54.8{\scriptsize$\pm$4.2} \\
272 & 60.7{\scriptsize$\pm$5.8} & 63.5{\scriptsize$\pm$4.6} & 75.4{\scriptsize$\pm$3.5} & 52.8{\scriptsize$\pm$4.8} & 33.5{\scriptsize$\pm$6.4} & 57.0{\scriptsize$\pm$2.4} \\
352 & 59.7{\scriptsize$\pm$5.6} & 61.4{\scriptsize$\pm$6.5} & 77.1{\scriptsize$\pm$4.0} & 48.6{\scriptsize$\pm$8.0} & 47.3{\scriptsize$\pm$4.8} & 56.1{\scriptsize$\pm$4.5} \\
\bottomrule
\end{tabular}}
}{\emph{(scaling table renders here as runs land)}}
\caption{\textbf{Model-size scaling.} Macro-F1 (mean$\pm$SD) across a width sweep
($d_{\mathrm{model}}\in\{96,\dots,352\}$) for Vanilla, MoR-general and bioMoR, averaged over
the $8$ single-cell and $3$ P-NET suites. bioMoR holds a flat, high accuracy across widths
while the vanilla stack---at $4\times$ the recursion-stack parameters---degrades.}
\label{fig:scaling_cv5}
\end{table}

\begin{table}[t]
\centering
\resizebox{\columnwidth}{!}{%
\begin{tabular}{lcccccc}
\toprule
Dataset & $M{=}128$ & $M{=}256$ & $M{=}512$ & $M{=}1024$ & $M{=}2048$ & Linear (all feat.) \\
\midrule
Baron & 79.1\,$\pm$\,2.8 & 78.1\,$\pm$\,1.4 & 79.6\,$\pm$\,4.4 & 82.6\,$\pm$\,2.0 & 80.4\,$\pm$\,2.1 & 91.1\,$\pm$\,2.0 \\
Lung & 79.1\,$\pm$\,1.2 & 80.2\,$\pm$\,0.8 & 79.8\,$\pm$\,1.9 & 81.3\,$\pm$\,2.0 & 81.1\,$\pm$\,1.9 & 74.8\,$\pm$\,0.2 \\
Muraro & 75.3\,$\pm$\,6.4 & 72.4\,$\pm$\,2.7 & 68.3\,$\pm$\,3.3 & 75.7\,$\pm$\,4.5 & 73.8\,$\pm$\,6.0 & 95.7\,$\pm$\,2.2 \\
Oeso. & 69.6\,$\pm$\,1.4 & 69.8\,$\pm$\,0.6 & 71.5\,$\pm$\,1.6 & 73.6\,$\pm$\,0.5 & 72.2\,$\pm$\,1.6 & 69.5\,$\pm$\,1.1 \\
Seger. & 71.4\,$\pm$\,5.3 & 70.3\,$\pm$\,4.1 & 73.8\,$\pm$\,5.4 & 69.4\,$\pm$\,14.7 & 68.3\,$\pm$\,3.3 & 89.5\,$\pm$\,5.0 \\
Spleen & 58.2\,$\pm$\,0.9 & 58.6\,$\pm$\,1.6 & 59.9\,$\pm$\,0.8 & 61.1\,$\pm$\,2.3 & 61.6\,$\pm$\,4.0 & 64.7\,$\pm$\,0.1 \\
T-cell & 68.9\,$\pm$\,1.0 & 69.9\,$\pm$\,0.3 & 71.5\,$\pm$\,0.4 & 71.8\,$\pm$\,2.3 & 73.0\,$\pm$\,0.3 & 78.1\,$\pm$\,0.3 \\
Xin & 64.7\,$\pm$\,9.4 & 64.8\,$\pm$\,6.0 & 63.7\,$\pm$\,10.3 & 64.1\,$\pm$\,0.6 & 62.1\,$\pm$\,2.8 & 97.4\,$\pm$\,0.6 \\
\midrule
\textbf{Mean macro-F1} & \textbf{70.8} & \textbf{70.5} & \textbf{71.0} & \textbf{72.5} & \textbf{71.6} & \textbf{82.6} \\
\bottomrule
\end{tabular}}
\caption{\textbf{Marker-budget headroom (3-seed held-out).} Single-cell macro-F1 as the
marker budget $M$ grows, against the full-feature linear baseline. Even when every feature
is a marker, a residual to the linear model remains, so the gap is architectural rather than
a token-budget limitation.}
\label{tab:mbudget}
\end{table}

\begin{table}[t]
\centering
\resizebox{\columnwidth}{!}{%
\begin{tabular}{lcccc}
\toprule
Dataset & Learned (random init) & $+$ bio warm-start & $+$ stronger init & $+$ annealed anchor \\
\midrule
Baron & 82.3\,$\pm$\,2.2 & 80.9\,$\pm$\,3.8 & 81.6\,$\pm$\,1.8 & 66.8\,$\pm$\,6.2 \\
Lung & 79.5\,$\pm$\,1.4 & 79.8\,$\pm$\,0.8 & 78.4\,$\pm$\,0.6 & 70.0\,$\pm$\,2.7 \\
Muraro & 74.5\,$\pm$\,5.4 & 75.3\,$\pm$\,4.9 & 76.6\,$\pm$\,5.4 & 75.9\,$\pm$\,9.3 \\
Oeso. & 69.3\,$\pm$\,0.7 & 71.1\,$\pm$\,0.9 & 70.2\,$\pm$\,1.0 & 57.4\,$\pm$\,2.6 \\
Seger. & 74.0\,$\pm$\,5.3 & 69.9\,$\pm$\,7.4 & 74.5\,$\pm$\,5.7 & 53.0\,$\pm$\,5.0 \\
Spleen & 59.0\,$\pm$\,0.2 & 59.1\,$\pm$\,0.7 & 58.7\,$\pm$\,1.2 & 52.5\,$\pm$\,1.2 \\
T-cell & 68.9\,$\pm$\,0.7 & 69.1\,$\pm$\,1.1 & 69.0\,$\pm$\,0.6 & 54.7\,$\pm$\,2.4 \\
Xin & 69.2\,$\pm$\,5.0 & 67.2\,$\pm$\,8.4 & 68.5\,$\pm$\,6.1 & 69.5\,$\pm$\,7.7 \\
\midrule
\textbf{Mean macro-F1} & \textbf{72.1} & \textbf{71.5} & \textbf{72.2} & \textbf{62.5} \\
\quad $\Delta$ vs.\ random & $+0.0$ & $-0.6$ & $+0.1$ & $-9.6$ \\
\bottomrule
\end{tabular}}
\caption{\textbf{Stress-testing the biological warm-start} (single-cell macro-F1,
mean$\pm$std). From a randomly-initialised learned graph we add a biological warm-start,
then a stronger init, then an annealed anchor $\lambda(t)\lVert A_{\text{learned}}-A_{\text{bio}}\rVert_F^2$
that keeps the graph near biology early. The $\Delta$ row is the mean gain over random
init; forcing biology in more strongly does not beat learning the graph from data.}
\label{tab:anchor}
\end{table}

\begin{table}[t]
\centering
\resizebox{\columnwidth}{!}{%
\begin{tabular}{lccc}
\toprule
Cohort & bioMoR (mut+CNV) & Geneformer & scGPT \\
\midrule
Prostate & 78.2\,$\pm$\,2.2 & 78.2\,$\pm$\,3.3 & 40.1\,$\pm$\,0.0 \\
BLCA & 40.9\,$\pm$\,1.9 & 47.7\,$\pm$\,3.7 & 40.4\,$\pm$\,0.0 \\
STAD & 52.2\,$\pm$\,6.8 & 46.1\,$\pm$\,5.5 & 35.7\,$\pm$\,0.0 \\
\midrule
\textbf{Mean} & \textbf{57.1} & \textbf{57.3} & \textbf{38.7} \\
\bottomrule
\end{tabular}}
\caption{\textbf{bioMoR vs.\ gene-vocabulary foundation models} on the P-NET cohorts
(macro-F1, mean$\pm$std over seeds). Geneformer (fine-tuned) and scGPT (frozen embedding
$+$ logistic probe) are mapped onto each cohort's HUGO gene symbols with a per-gene
mutation/copy-number alteration burden, on the same splits as bioMoR. Bulk DNA-alteration
input is out-of-distribution for these single-cell-RNA models.}
\label{tab:fm}
\end{table}

% Calibration / uncertainty table removed.

\begin{table}[t]
\centering
\resizebox{0.85\columnwidth}{!}{%
\begin{tabular}{lccc}
\toprule
Depth $K$ & Shared (ours) & Independent & Reduction \\
\midrule
1 & 74{,}880 & 74{,}880 & 1$\times$ \\
2 & 74{,}880 & 149{,}760 & 2$\times$ \\
3 & 74{,}880 & 224{,}640 & 3$\times$ \\
4 & 74{,}880 & 299{,}520 & 4$\times$ \\
6 & 74{,}880 & 449{,}280 & 6$\times$ \\
8 & 74{,}880 & 599{,}040 & 8$\times$ \\
\bottomrule
\end{tabular}}
\caption{\textbf{Parameter reduction.} One shared block versus $K$ independent layers at
matched width: an exact $K\times$ reduction, present before any training.}
\label{tab:param}
\end{table}

\begin{table}[t]
\centering
\resizebox{\columnwidth}{!}{%
\begin{tabular}{lccccc}
\toprule
Dataset & $M{=}16$ & $M{=}32$ & $M{=}64$ & $M{=}128$ & $M{=}256$ \\
\midrule
\multicolumn{6}{l}{\emph{Single-cell (genomap)}} \\
\quad Baron & 33.7\,$\pm$\,7.1 & 38.5\,$\pm$\,2.1 & 53.9\,$\pm$\,1.9 & 62.6\,$\pm$\,5.1 & 66.9\,$\pm$\,3.4 \\
\quad Lung & 30.8\,$\pm$\,5.9 & 38.1\,$\pm$\,1.7 & 61.5\,$\pm$\,1.9 & 71.4\,$\pm$\,1.1 & 78.4\,$\pm$\,0.8 \\
\quad Muraro & 55.8\,$\pm$\,3.6 & 59.4\,$\pm$\,1.3 & 67.9\,$\pm$\,5.6 & 68.5\,$\pm$\,4.1 & 76.2\,$\pm$\,5.2 \\
\quad Oeso. & 23.1\,$\pm$\,1.3 & 30.2\,$\pm$\,3.9 & 42.0\,$\pm$\,2.4 & 53.9\,$\pm$\,4.4 & 63.2\,$\pm$\,2.4 \\
\quad Seger. & 41.2\,$\pm$\,12.1 & 45.3\,$\pm$\,5.6 & 59.3\,$\pm$\,4.0 & 59.1\,$\pm$\,1.6 & 55.4\,$\pm$\,7.0 \\
\quad Spleen & 18.7\,$\pm$\,2.4 & 26.8\,$\pm$\,3.0 & 38.1\,$\pm$\,2.3 & 48.1\,$\pm$\,5.0 & 59.4\,$\pm$\,0.9 \\
\quad T-cell & 25.5\,$\pm$\,3.2 & 36.6\,$\pm$\,5.4 & 43.1\,$\pm$\,5.2 & 47.4\,$\pm$\,2.2 & 61.5\,$\pm$\,1.7 \\
\quad Xin & 50.5\,$\pm$\,5.9 & 46.2\,$\pm$\,1.6 & 60.8\,$\pm$\,5.6 & 64.1\,$\pm$\,3.3 & 72.4\,$\pm$\,11.7 \\
\midrule
\multicolumn{6}{l}{\emph{Multi-omics (Reactome/P-NET)}} \\
\quad Prostate & 74.3\,$\pm$\,3.0 & 74.7\,$\pm$\,2.7 & 78.0\,$\pm$\,3.0 & 76.7\,$\pm$\,2.6 & 77.3\,$\pm$\,0.2 \\
\quad BLCA & 49.8\,$\pm$\,2.2 & 46.0\,$\pm$\,3.1 & 49.3\,$\pm$\,2.5 & 48.1\,$\pm$\,6.8 & 55.7\,$\pm$\,3.5 \\
\quad STAD & 54.2\,$\pm$\,3.1 & 51.0\,$\pm$\,2.2 & 44.0\,$\pm$\,5.6 & 51.2\,$\pm$\,2.8 & 49.4\,$\pm$\,4.1 \\
\midrule
\textbf{Mean macro-F1} & \textbf{41.6} & \textbf{44.8} & \textbf{54.4} & \textbf{59.2} & \textbf{65.1} \\
\bottomrule
\end{tabular}}
\caption{\textbf{Marker-token budget.} Macro-F1 (mean$\pm$std over 3 seeds) as
the marker budget $M$ shrinks from 256 to 16. A few dozen to a few hundred tokens recover
most of the full-gene accuracy. Bottom row: mean over all 11 datasets.}
\label{tab:token}
\end{table}



\section{Dataset Details}
\label{app:data}
We use $8$ genomap single-cell datasets \cite{islam2023cartography} converted to
plain CSV (expression + labels + stratified split) -- Baron, Lung, Muraro, Oesophagus,
Segerstolpe, Spleen, T-cell and Xin -- $3$ Reactome/P-NET multi-omics
cohorts (prostate, BLCA, STAD) whose fixed Reactome pathway tokens pool mutation,
copy-number and expression channels. Per-dataset sample, feature and class counts are
read directly from the result files and summarised in Table~\ref{tab:datasets}.

\begin{table}[h]
\centering
\resizebox{\columnwidth}{!}{%
\begin{tabular}{llrrr}
\toprule
Dataset & Modality & Samples & Features & Classes \\
\midrule
Baron & single-cell & 8{,}569 & 2{,}000 & 13 \\
Lung & single-cell & 57{,}020 & 1{,}089 & 25 \\
Muraro & single-cell & 2{,}122 & 2{,}000 & 9 \\
Oeso. & single-cell & 87{,}947 & 1{,}089 & 18 \\
Seger. & single-cell & 2{,}127 & 2{,}000 & 11 \\
Spleen & single-cell & 94{,}129 & 1{,}089 & 28 \\
T-cell & single-cell & 24{,}007 & 1{,}089 & 7 \\
Xin & single-cell & 1{,}492 & 2{,}000 & 4 \\
\midrule
Prostate & multi-omics & 1{,}011 & 1{,}223 & 2 \\
BLCA & multi-omics & 404 & 1{,}258 & 2 \\
STAD & multi-omics & 414 & 1{,}258 & 2 \\
PM (\texttt{pan\_meta}, mut+CNV) & multi-omics & 8{,}893 & 10{,}333 & 2 \\
PC (\texttt{pan\_meta}, expr) & multi-omics & 8{,}586 & 10{,}013 & 2 \\
\bottomrule
\end{tabular}}
\caption{The 13 datasets used throughout: $8$ genomap single-cell suites
$3$ Reactome/P-NET multi-omics cohorts. Counts are read directly from the
result files.}
\label{tab:datasets}
\end{table}

\section{Effective-FLOPs Accounting}
\label{app:flops}
The compute numbers report per-sample FLOPs of the recursive transformer stack, the
only component routing changes. One application of the shared block to $a$ tokens costs
$\phi(a)=4a^2 d + 4\,a\,d\,d_{\mathrm{ff}}$; the nominal fixed-depth cost is
$\Phi_{\mathrm{nom}}=K\,\phi(M)$ and the effective cost sums one block over the tokens
active at each step, $\Phi_{\mathrm{eff}}=\sum_{t=1}^{K}\phi(a_t)$, where $a_t$ is the
mean number of markers the expert-choice funnel keeps at step $t$. Gene embedding and
marker selection are $\mathcal{O}(Nd)$, identical across routing modes, and excluded
from this stack-level comparison.

\paragraph{End-to-end accounting.}
For completeness we account for the whole forward pass, not only the stack. Three parts
scale with the full gene count $N$: gene embedding $\Theta(Nd)$; the cross-attention
marker router, whose $M$ queries attend over all $N$ gene keys at
$\Theta(MNd)$; and the optional gene-graph smoother, which at rank $r$ costs
$\Theta(Nr)$ (never the $N^2$ dense form). Everything after marker selection -- the
recursive stack, pooling and head -- scales with the compressed budget $M\ll N$ at
$\Theta(M^2d)$ per pass. So the router's $\Theta(MNd)$ term dominates the $N$-scaling
and is \emph{linear} in $N$ (versus the $\Theta(N^2 d)$ self-attention a full-gene
transformer would pay), while the quadratic cost is paid only on the $M$ markers. The
architecture ablations of the main-paper efficiency ladder vary only the stack, so the stack-level
ratios there are the correct comparison for the routing/sharing claims; the router and
embedding terms are shared by every variant. Measured wall-clock and peak-memory
profiling on matched hardware is a straightforward addition we leave to the camera-ready.

\section{Theoretical Foundation of the Router}
\label{app:theory}
This appendix gives the mathematical and biological grounding for bioMoR's
biology-informed router.

\paragraph{Routing as conditional computation.}
The router implements \emph{conditional computation}: a learned policy routes each token
to a token-specific amount of compute, the gating principle of sparsely-gated mixtures
of experts \cite{shazeer2017outrageously}, reused by Mixture-of-Depths
\cite{raposo2024mixture} and Mixture-of-Recursions \cite{bae2025mixture}; the
``experts'' here are recursion \emph{depths} of one shared block, which couples adaptive
computation \cite{graves2016adaptive} to weight sharing.

\paragraph{The differentiable handle.}
The discrete top-$k$ has zero gradient almost everywhere, so bioMoR keeps the soft
probability of the chosen route as a multiplicative \emph{gate} on the block output,
$\mathbf{o}_m=g_m\,f_\theta(\mathbf{h}_m)$ with $g_m=\sigma(\tilde r_m)$. Because $g_m$
is smooth in $\mathbf{w}_r$, the chain
$\mathcal{L}\!\leftarrow\!\mathbf{o}_m\!\leftarrow\!g_m\!\leftarrow\!\mathbf{w}_r$ is
unbroken: the hard choice routes, the soft gate carries the gradient. The biological
prior of the routing logit of the main paper (Eq.~1) is an additive constant in this logit, so it shifts
the decision without breaking this path.

\paragraph{Biological foundation of the prior.}
Two facts from single-cell biology motivate the prior. A small set of \emph{marker}
genes carries most discriminative signal
\cite{ianevski2022fully,franzen2019panglaodb,hu2023cellmarker}, so compute should be
allocated unevenly. And gene co-expression networks are approximately scale-free: a few
high-degree \emph{hub} genes (master regulators) exert outsized influence. We
operationalise ``hub'' as eigenvector centrality on the genomap co-expression graph
$\mathbf{W}$: the leading eigenvector $\mathbf{v}$ of
$\mathbf{W}\mathbf{v}=\lambda\mathbf{v}$ scores each gene by the centrality of its
neighbours, recursively, and we $z$-score it to form $\pi$.

\paragraph{Empirical-Bayes reading.}
The routing logit of the main paper (Eq.~1) is a log-linear prior on the routing decision, with
$\beta_t\,\pi_m$ a Gaussian-like prior mean and $\beta_t=\beta_0(1-\text{progress})$ a
shrinkage strength that decays as data evidence accumulates: prior-dominated when the
likelihood is uninformative (early training), data-dominated later.

\paragraph{Leakage safety.}
$\mathbf{W}$ is computed from training-split \emph{expression only}; no cell-type label
enters it. The prior therefore injects network topology, not the answer, which is why a
gene-discovery claim under this prior is not circular, unlike a prior built from curated
cell-type marker lists.

\paragraph{Optional pathway-graph smoothing.}
The additive bias treats genes independently. Co-pathway genes should route coherently,
which one obtains by smoothing the logits over $\mathbf{W}$ with the normalised
Laplacian $\mathbf{L}=\mathbf{I}-\mathbf{D}^{-1/2}\mathbf{W}\mathbf{D}^{-1/2}$,
$\hat{\mathbf{r}}^{(t)}=\tilde{\mathbf{r}}^{(t)}-\gamma\,\mathbf{L}\tilde{\mathbf{r}}^{(t)}$
(graph-Laplacian regularisation): a gene borrows routing strength from its network
neighbours. This is one sparse matrix-vector product per step and adds no transformer
parameters; we expose it as an option and leave its evaluation to future work.

\bibliographystyle{aaai}
\bibliography{refs}
\end{document}
